\begin{frame}[fragile]
\frametitle{Q4(1): 問題設定}

\begin{alertblock}{問題}
$REGULAR_{\text{TM}}$ および $\overline{REGULAR_{\text{TM}}}$ が、どちらもチューリング認識不可能（Turing-unrecognizable）であることを示せ。
\end{alertblock}

\textbf{証明のポイント:}
\begin{itemize}
    \item 既知の事実：$\overline{A_{\text{TM}}}$ はチューリング認識不可能である。
    \item 示すべきこと：
    \begin{enumerate}
        \item $\overline{A_{\text{TM}}} \le_{m} REGULAR_{\text{TM}}$ を示す $\Rightarrow REGULAR_{\text{TM}}$ が認識不可能。
        \item $\overline{A_{\text{TM}}} \le_{m} \overline{REGULAR_{\text{TM}}}$ を示す $\Rightarrow \overline{REGULAR_{\text{TM}}}$ が認識不可能。
    \end{enumerate}
    \item 性質：$L_1 \le_m L_2 \iff \overline{L_1} \le_m \overline{L_2}$ を利用する。
\end{itemize}
\end{frame}

\begin{frame}[fragile]
\frametitle{Q4(1) 第1部:  $\overline{REGULAR_{\text{TM}}}$  の認識不可能性}

$A_{\text{TM}} \le_m \overline{REGULAR_{\text{TM}}}$ を示すため、マシン $G$ を構成する。
$G = $ 「入力 $\langle M, w \rangle$ に対し、以下の TM $M'$ を構成して出力する」
\begin{block}{$M'$（入力 $x$）}
    \begin{enumerate}
        \item $x \neq 0^n 1^n$ ならば \textbf{拒否}する。
        \item $x = 0^n 1^n$ ならば $M$ を $w$ に対して実行する。
        \item $M$ が受理すれば $M'$ も \textbf{受理}する。
    \end{enumerate}
\end{block}

\textbf{結果の確認:}
\begin{itemize}
    \item $M$ が $w$ を受理する $\Rightarrow L(M') = \{0^n 1^n\}$（非正規）$\Rightarrow \langle M' \rangle \in \overline{REGULAR_{\text{TM}}}$
    \item $M$ が $w$ を受理しない $\Rightarrow L(M') = \emptyset$（正規）$\Rightarrow \langle M' \rangle \notin \overline{REGULAR_{\text{TM}}}$
\end{itemize}
よって $A_{\text{TM}} \le_m \overline{REGULAR_{\text{TM}}}$。これより $\overline{A_{\text{TM}}} \le_m REGULAR_{\text{TM}}$ となり、$REGULAR_{\text{TM}}$ は認識不可能。
\end{frame}

\begin{frame}[fragile]
\frametitle{Q4(1) 第2部: $REGULAR_{\text{TM}}$ の認識不可能性}

$A_{\text{TM}} \le_m REGULAR_{\text{TM}}$ を示すため、マシン $F$ を構成する。
$F = $ 「入力 $\langle M, w \rangle$ に対し、以下の TM $M_2$ を構成して出力する」
\begin{block}{$M_2$（入力 $x$）}
    \begin{enumerate}
        \item $x = 0^n 1^n$ ならば \textbf{受理}する。
        \item $x \neq 0^n 1^n$ ならば $M$ を $w$ に対して実行する。
        \item $M$ が受理すれば $M_2$ も \textbf{受理}する。
    \end{enumerate}
\end{block}

\textbf{結果の確認:}
\begin{itemize}
    \item $M$ が $w$ を受理する $\Rightarrow L(M_2) = \Sigma^*$（正規）$\Rightarrow \langle M_2 \rangle \in REGULAR_{\text{TM}}$
    \item $M$ が $w$ を受理しない $\Rightarrow L(M_2) = \{0^n 1^n\}$（非正規）$\Rightarrow \langle M_2 \rangle \notin REGULAR_{\text{TM}}$
\end{itemize}
よって $A_{\text{TM}} \le_m REGULAR_{\text{TM}}$。これより $\overline{A_{\text{TM}}} \le_m \overline{REGULAR_{\text{TM}}}$ となり、$\overline{REGULAR_{\text{TM}}}$ は認識不可能。
\end{frame}

\begin{frame}[fragile]
\frametitle{Q4(1): 結論}

\begin{itemize}
    \item \textbf{前半:} $\overline{A_{\text{TM}}} \le_m REGULAR_{\text{TM}}$ を示した。\\
    $\overline{A_{\text{TM}}}$ は認識不可能であるため、写像還元の性質により $REGULAR_{\text{TM}}$ もチューリング認識不可能である。
    \vspace{1em}
    \item \textbf{後半:} $\overline{A_{\text{TM}}} \le_m \overline{REGULAR_{\text{TM}}}$ を示した。\\
    同様に、$\overline{REGULAR_{\text{TM}}}$ もチューリング認識不可能である。
\end{itemize}

\begin{center}
    \textbf{以上より、$REGULAR_{\text{TM}}$ と その補集合はどちらも認識不可能である。}
\end{center}
\end{frame}